\section{Experiments}

\subsection{Setup}

\paragraph{Models and agent scaffold.}
We use Claude 4 Sonnet, Claude 4.5 Sonnet, Claude 4.5 Opus, and GPT-5.2 as the primary backbone LLMs to ensure comparability with published baselines. We use SWE-Agent \citep{sweagent} as the baseline scaffold. Our CCA agent replaces the SWE-Agent stack with the Confucius Code Agent.
We also report results from the Live-SWE-Agent \citep{xia2025live} as baseline, while keeping the tool environment and repository setup identical. 

\paragraph{Benchmark.}
For main results, we evaluate CCA on the SWE-Bench-Pro \citep{sbp} public split consisting of 731 tasks, following the identical environment configuration and infrastructure used by the SWE-Agent baseline \citep{sweagent}. We also report results from the SWE-Bench-Verified \citep{swebench} consisting of 500 tasks to compare with existing open-sourced coding agents, including SWE-Agent and OpenHands \citep{openhands}.

\paragraph{Metrics.}
We follow the official SWE-Bench-Pro Resolve Rate metric, defined as the percentage of tasks for which the agent’s proposed patch successfully passes all repository-provided tests without human intervention.\footnote{\url{https://scale.com/leaderboard/swe_bench_pro_public}} Each trial is repeated with different random seeds for trajectory sampling to account for randomness in tool invocation and LLM responses. We report mean Resolve Rate (Pass@1) across three runs.

\subsection{Main Results on SWE-Bench-Pro}

\begin{table}[h]
\centering
\begin{tabular}{l l c}
\toprule
\textbf{Backbone Model} & \textbf{Scaffold} & \textbf{Resolve Rate (Pass@1)} \\
\midrule

\multirow{2}{*}{Claude 4 Sonnet}
    & SWE-Agent \citep{sweagent} & 42.7 \\
    & CCA           & 45.5 \\
\midrule

\multirow{3}{*}{Claude 4.5 Sonnet}
    & SWE-Agent  & 43.6 \\
    & Live-SWE-Agent \citep{xia2025live}       & 45.8 \\
    & CCA            & 52.7 \\
\midrule

\multirow{2}{*}{Claude 4.5 Opus}
    & Anthropic*  & 52.0 \\
    & CCA            & \textbf{54.3} \\
\midrule

\multirow{2}{*}{GPT-5.2}
    & OpenAI* & 56.0 \\
    & CCA     & \textbf{59.0} \\
\bottomrule
\end{tabular}
\caption{SWE-Bench-Pro public split comparison across scaffolds and backbone models.
All methods share identical environments; improvements arise solely from the agent scaffolds.
(* Anthropic's proprietary scaffold, from Claude Opus 4.5 system card; * OpenAI's proprietary scaffold, from GPT-5.2 report.)}
\label{tab:swepro}
\end{table}

\begin{table}[t]
\centering
\begin{tabular}{l l l c}
\toprule
\textbf{Backbone Model} & \textbf{Context Management} & \textbf{Tool Use} & \textbf{Resolve Rate (Pass@1)} \\
\midrule

\multirow{2}{*}{Claude 4 Sonnet}
    & No  & advanced & 42.0 \\
    & Yes & advanced & 48.6 \\
\midrule

\multirow{3}{*}{Claude 4.5 Sonnet}
    & No  & simple   & 44.0 \\
    & No  & advanced & 51.0 \\
    & Yes & advanced & \textbf{51.6} \\
\bottomrule
\end{tabular}
\caption{Ablation on context management (hierarchichal working memory \& context compression) and tool-use sophistication. Results are obtained from evaluating CCA on a 100-example subset of the SWE-Bench-Pro public set.}
\label{tab:context_mgmt}
\end{table}


\Cref{tab:swepro} summarizes our main results on the SWE-Bench-Pro public split. Under identical environment and tool conditions, CCA consistently surpasses the SWE-Agent baseline across settings with different backbone models. 
With Claude 4 Sonnet, CCA reaches Resolve@1 at \textbf{45.5\%}. With Claude 4.5 Sonnet, CCA reaches \textbf{52.7\%}, largely surpassing the best research-grade coding agent, Live-SWE-Agent, at 45.8\%.
With Claude 4.5 Opus, CCA achieves \textbf{54.3\%}, achieving higher performance than results reported by Anthropic. 
And with GPT-5.2, we achieve \textbf{59.0\%}, surpassing the official result reported by OpenAI and setting new leading performance on SWE-Bench-Pro.
These improvements arise purely from stronger agentic scaffolding, i.e. enhanced orchestration, context management, and tool-use extensions, rather than differences in backbone models or evaluation setups.
More broadly, these results underscore the central role of scaffolding: even a weaker model equipped with a strong agent scaffold (Claude 4.5 Sonnet + CCA at \textbf{52.7\%}) can outperform a stronger model (Claude 4.5 Opus + Anthropic's proprietary scaffold at \textbf{52.0\%}). 

Because Claude Code (CC) does not expose a programmatic tool interface compatible with containerized evaluation environments such as SWE-rex \citep{SWEReX2025}, we cannot compare CCA with CC's results on SWE-Bench-Pro. Instead, to provide a qualitative comparison, we constructed a small curated benchmark (a mini PyTorch-Bench) and executed CC solutions using the Claude Code CLI directly on a host machine where CC is installed w/o a Docker-based runtime, as seen in Appendix~\ref{sec:case_studies}. These complementary experiments highlight behavioral differences between CC and CCA in realistic debugging and development tasks, but they are not directly comparable to SWE-Bench-Pro due to differences in execution environments.


\subsection{Meta-Agent Learned Tool-Use}
\label{sec:meta-agent-tool-use}

CCA’s tool-use behavior is not purely hand-engineered; instead, it is \emph{learned} through the Meta-agent, which automatically refines how the agent invokes tools such as file editors and command-line utilities. 
To measure the contribution of this learned tool-use stack, we perform an ablation that disables these Meta-agent–derived tools and instead reverts CCA to a simpler,    ``naive'' tool-use pattern similar to traditional SWE-Agent-like scaffolds with simple file editing and command-line operations only.
\Cref{tab:context_mgmt} reports the results of this ablation alongside a separate ablation on context management. Experiments are conducted on a 100-example subset of the SWE-Bench-Pro public set. As shown in the Claude 4.5 Sonnet rows, removing the learned tool-use features leads to a large decline in Resolve@1, even when the context management method is held constant. This confirms that tool-use conventions learned by the Meta-agent are a major driver of CCA’s performance, independent of (and complementary to) hierarchical working memory and context compression.

The function below shows an example of meta-agent's improvement on CCA's prompts. The part emphasized here is a {failure message} that gets raised when the best match is not exact. This message was iteratively refined by the meta-agent to be maximally actionable for an LLM: it shows the closest match, provides a minimal unified diff, and gives explicit constraints that force the agent to keep using \texttt{<file\_edit>} and to repair the \texttt{<find>/<find\_after>} tag into an exact \texttt{<line\_number>|<exact\_line\_content>} format rather than switching tools or attempting unsafe workarounds.

\begin{lstlisting}[language=Python,caption={Excerpt of the file-edit chunk matcher; non-prompt logic is collapsed as "\texttt{...}", while the meta-agent-refined prompt is shown verbatim.},label={lst:file_edit_prompt_excerpt}]
def _get_matched_chunk(...):
    # ... 
    
    if similarity == 1.0: 
        return matched_chunk

    # ... 

    # --- meta-agent-refined prompt begins
    raise ValueError(
        dedent(
            """\
        No exact occurrence found for the search string you provided.
        
        # ... shows match and diff
        
        ACTION REQUIRED: Please update your `<find>` or `<find_after>` tag to 
        match the exact content in the file with `<line_number>|
        <exact_line_content>` format.
        
        IMPORTANT: YOU MUST CONTINUE USING <file_edit> tag until successful. 
        DO NOT attempt alternative approaches such as:
        - Creating a new file to override the existing one
        - Using command line tools (e.g., `sed`, `awk`, etc.)
        
        Continue refining your `<find>` or `<find_after>` tag until it exactly 
        matches the file content.
        """
        ).format(...)
    )
    # --- meta-agent-refined prompt ends
\end{lstlisting}


\subsection{Evaluations on Long-context Reasoning}

\subsubsection{Context Management}

To quantify the impact of hierarchical working memory and context compression, we evaluate CCA on a subset of SWE-Bench-Pro, where both variants (with and without context management) successfully produced executable solutions.\footnote{Without any context control, many trajectories exceed model token limits and fail to complete, hence the restricted subset.} 
Results in \Cref{tab:context_mgmt} demonstrate a clear improvement in problem resolution when hierarchical memory and context compression are enabled. For Claude 4 Sonnet, advanced context management improves Resolve@1 from $42.0$ to $48.6$ on this subset (a +6.6 performance gain). On Claude 4.5, the improvement between the no-context-management and advanced variants is smaller, but both substantially outperform the simple tool-use configuration.  
This supports the hypothesis that structured context compression not only prevents overflow but also improves reasoning quality by enforcing periodic consolidation of long-horizon plans.
To ensure a controlled comparison, we use the same backbone LLM for both context compression and the main orchestrator agent across all experiments; see detailed study in the Appendix~\ref{app:context}. 







\subsubsection{Endless-Read Robustness}

We further analyze CCA's robustness under tasks that require editing multiple files. Each SWE-Bench-Pro task is grouped by the number of modified files (``edited-file bucket''), and we measure the Resolve Rate within each group. As shown in \Cref{tab:endless_read}, the agent maintains stable performance across varying edit volumes, with only moderate regression when more files are touched. The degradations likely stem from cumulative localization uncertainty and compounding diffs, suggesting future work on finer-grained diff validation and multi-file dependency tracking.
Overall, these results show that CCA’s hierarchical memory and context compression yield substantial gains in both efficiency and robustness for long-context reasoning.

\begin{table}[h]
\centering
\begin{tabular}{l c c}
\toprule
\textbf{Edited Files} & \textbf{Resolve Rate (Pass@1)} & \textbf{\# Samples} \\
\midrule
1--2 files  & 57.8 & 294 \\
3--4 files  & 49.2 & 203 \\
5--6 files  & 44.1 & 86  \\
7--10 files & 52.6 & 38  \\
10+ files   & 44.4 & 18  \\
\bottomrule
\end{tabular}
\caption{CCA's resolve rate on SWE-Bench-Pro as a function of the number of files modified. 
Performance remains robust even for multi-file refactoring scenarios.}
\label{tab:endless_read}
\end{table}


\subsection{Evaluations on Long-term Memory}

We next study CCA’s \emph{note-taking} module, designed to accumulate durable cross-session memory. Unlike transient hierarchical working memory, the note-taking agent asynchronously summarizes each session into structured Markdown notes with multiple steps of reasoning, which capture both successful strategies and failure cases. This persistent ``memory'' is then available for retrieval in subsequent tasks, supporting test-time self-improvement.

Since no public benchmark explicitly evaluates memory in coding agents, we assess CCA’s memory module by running it on two consecutive passes, i.e., with memory maintained, of SWE-Bench-Pro instances. During the first run, the note-taking agent analyzes each trajectory and produces persistent notes for 151 instances—skipping cases where no meaningful insight can be distilled. We then rerun exactly these 151 tasks, providing CCA with the corresponding note directory to measure how prior experience improves efficiency and solution quality.
For Run 1: We execute the task from scratch (no context editing either); use note taker agent to write  down notes. For Run 2: We pass the notes from Run~1 to CCA and rerun.

\begin{table}[h]
\centering
\small
\begin{tabular}{l l l l}
\toprule
\textbf{Trial} & \textbf{Avg.\ Turns ($\downarrow$)} & \textbf{Avg.\ Token Cost ($\downarrow$)} & \textbf{Resolve Rate (Pass@1, $\uparrow$)} \\
\midrule
Run 1 & 64 & 104k & 53.0 \\
Run 2 & 61 \textcolor{red}{\textbf{(-3)}} & 93k \textcolor{red}{\textbf{(-11k)}} & 54.4 \textcolor{red}{\textbf{(+1.4)}} \\
\bottomrule
\end{tabular}
\caption{CCA performance across repeated runs using notes. For Run 1, all tasks are processed from scratch, and notes are taken and stored. For Run 2, each task resolving session is accompanied by the taken notes as the long-term memory. Token cost excludes system prompt tokens; the underlying model is Claude 4.5 Sonnet.}
\label{tab:note_ablation}
\end{table}

Cumulative note-taking reduces the iteration turns (from 64 to 61) and the token cost (from $104$k to $93$k), and also yield improvements on resolve rate (from 53\% to 54.4\%). These gains indicate that the notes distilled in the first run capture actionable, reusable knowledge. In effect, the note-taking system provides CCA with a lightweight form of \emph{cross-session learning}, enabling more efficient reasoning and more reliable patch generation in subsequent attempts. A detailed example of the notes produced by the note-taking agent is provided in Appendix~\ref{app:notes}. 


\subsection{Comparison with Open-Sourced Scaffolds on SWE-Bench-Verified}

We further conduct evaluations on the SWE-Bench-Verified benchmark \citep{swebench} to compare CCA against existing open-source scaffolds.\footnote{As of Dec 2025, OpenHands remains the strongest open-sourced coding agent on SWE-Bench-Verified, reported from SWE-Bench's official leaderboard.} Using Claude 4 Sonnet, CCA achieves a Resolve Rate of \textbf{74.6\%}, exceeding the strongest open-source system (OpenHands) under identical backbone conditions and outperforming a mini-SWE-Agent variant that relies on the more capable Claude 4.5 Sonnet model. These results reinforce the central role of agentic scaffolding: improved orchestration, memory handling, and tool-use abstractions can close—or even surpass—the gap introduced by differences in backbone model capability. We also observe that SWE-Bench-Verified is sensitive to Claude’s internal thinking budget; a detailed analysis appears in Appendix~\ref{app:thinking}.


\begin{table}[h]
\centering
\begin{tabular}{l l c}
\toprule
\textbf{Backbone Model} & \textbf{Scaffold} & \textbf{Resolve Rate (Pass@1)} \\
\midrule

\multirow{3}{*}{Claude 4 Sonnet}
    & SWE-Agent            & 66.6 \\
    & OpenHands \citep{openhands}            & 72.8 \\
    & \textbf{CCA}         & \textbf{74.6} \\
\midrule

Claude 4.5 Sonnet
    & mini-SWE-Agent       & 70.6 \\
\bottomrule
\end{tabular}
\caption{CCA performance on SWE-Bench-Verified. CCA matches the best open-source framework (OpenHands) under the same Claude 4 Sonnet backbone, and outperforms a mini-SWE-Agent variant even when that variant uses a stronger Claude~4.5~Sonnet backbone.}
\label{tab:swe_verified}
\end{table}



